\partstart{Conservation of Information}
\section{Thalean Conservation of Information}\label{sec:information}
\begin{principle}[title=Thalean conservation of information]
Admissible transduction may redistribute relational information between propagation and receipt, but may not erase the distinction structure required to reconstruct the completed handoff.
\end{principle}
This is the governing formal principle of QR. In this manuscript it means complete recoverability on the declared admitted domain. It does not assert that every scalar readout is conserved or equate the principle with an unproved physical information law.

\setcounter{theorem}{0}
\begin{proposition}[Classical information preservation]
For a map $\mathsf T:X\to Y$, a reconstruction map $\mathsf R:\im\mathsf T\to X$ satisfying $\mathsf R\mathsf T=\id_X$ exists if and only if $\mathsf T$ is injective.
\end{proposition}
\begin{proof}
If $\mathsf T(x_1)=\mathsf T(x_2)$, applying $\mathsf R$ gives $x_1=x_2$. Conversely, if $\mathsf T$ is injective, each point of its image has exactly one preimage; assigning that preimage defines $\mathsf R$.
\end{proof}

\setcounter{theorem}{1}
\begin{proposition}[Compositional preservation]
A composition of admitted injective event maps is injective, provided the domains match and later maps include rather than silently discard the retained outputs required by earlier maps.
\end{proposition}
\begin{proof}
For two events with reconstructors $\mathsf R_1,\mathsf R_2$, the reverse-order composition $\mathsf R_1\circ\mathsf R_2$ reconstructs $\mathsf T_2\circ\mathsf T_1$. Induction proves the claim for a finite history. The domain qualification ensures that the second map acts on the complete intermediate state.
\end{proof}

Fix a propagated map $p:X\to P$ on a finite classical domain. A completing receipt $r:X\to R$ makes $x\mapsto(p(x),r(x))$ injective.
\setcounter{theorem}{2}
\begin{proposition}[Receipt-capacity bound]
A completing receipt alphabet exists with size $k$ if and only if
% Original equation number(s) 14; expression unchanged.
\setcounter{equation}{13}
\begin{equation}
k\geq\max_{y\in p(X)}|p^{-1}(y)|.
\label{eq:receiptcapacity}
\end{equation}

\end{proposition}
\begin{proof}
Within one propagated fiber, receipt labels must all differ, proving necessity. Conversely, enumerate each fiber independently using labels $1,\ldots,k$. The pair of its propagated value and its fiber label identifies every input, proving sufficiency.
\end{proof}

Thus an injective propagated map requires only one receipt value. When a map identifies several inputs, the receipt must distinguish them within each propagated fiber.

For a deterministic injective encoding of a specified classical distribution,
% Original equation number(s) 15; expression unchanged.
\setcounter{equation}{14}
\begin{equation}
H(X)=H(P,R)=H(P)+H(R)-I(P{:}R).
\label{eq:shannon}
\end{equation}

and $H(R\mid P)=H(X\mid P)$. The mutual-information term removes double counting of correlations \citep{shannon1948}. For an isometric quantum encoding $\rho_{PR}=V\rho_XV^\dagger$, spectrum preservation gives
% Original equation number(s) 16; expression unchanged.
\setcounter{equation}{15}
\begin{equation}
S(\rho_X)=S(\rho_{PR})
=S(\rho_P)+S(\rho_R)-I(P{:}R).
\end{equation}

These are joint-output identities; neither marginal need be reversible. Logical recovery and thermodynamic reversibility remain different requirements \citep{bennett1973,watrous2018}.

\partstart{Scalar Complexity, Phase, and Mode}
\begingroup
\renewcommand{\thesection}{11A}
\section{Scalar Relay as Relational Accounting}\label{sec:scalarrelay}
\status{working relay definition; conditional algebra proved here; a native local realization remains to be supplied}
Fix a registered-object domain, its complexity functional, and an admitted update $R_i\xrightarrow{\mathsf T_i}R_{i+1}$. The objects refer to a consistent accounting boundary; when a propagated subsystem is used, the complementary receipt still belongs to the complete event. Define
\begin{equation}
\begin{aligned}
\Delta\mathcal C_i&:=\mathcal C(R_{i+1})-\mathcal C(R_i),\\
\mathcal C_{i+1}&=\mathcal C_i+\Delta\mathcal C_i,
& S_{i+1}&=S_iG^{\Delta\mathcal C_i}.
\end{aligned}
\tag{RC3}\label{eq:scalarrelay}
\end{equation}
The second line follows from the first and from (RC2). It is not yet a mechanism for computing the increment locally.

\begin{definition}[Scalar relay]
A Thalean transducer relays scalar complexity by admitting a relational transformation whose native complexity increment updates the registered scalar. The scalar need not be transported as an independent payload; it is reconstructed from continuity of relational accounting, using an admitted reference value or sufficient relational state.
\end{definition}
This is an implementation requirement on a relational interface, not a claim of communication without information. A downstream transducer must have access to the distinctions needed for its scalar readout. Global recoverability of the joint output does not guarantee that the propagated component alone supplies them.

\begin{proposition}[Compositional scalar accounting]\label{prop:scalarcomposition}
For a composable finite sequence on which one fixed $\mathcal C$ is defined,
\begin{equation}
\sum_{i=0}^{n-1}\Delta\mathcal C_i=\mathcal C(R_n)-\mathcal C(R_0),
\qquad
S_n=S_0\prod_{i=0}^{n-1}G^{\Delta\mathcal C_i}.
\tag{RC4}\label{eq:scalartelescope}
\end{equation}
Thus a known initial value and exact increments determine the final scalar. On a closed complete-state sequence $R_n=R_0$, the total increment is zero.
\end{proposition}
\begin{proof}
The sum telescopes. Exponentiating the result in base $G$ gives the product identity. Closure gives equal endpoint values. A returning projected readout need not be closure of the complete registered object.
\end{proof}

\subsection*{What the next transducer must receive}
Let $z$ be the complete admitted input to a fixed update $\mathsf T$, and let $\eta(z)$ denote exactly the relational data and prior local context available to the receiving interface. For brevity write
\[
\delta(z)=\mathcal C(\mathsf Tz)-\mathcal C(z).
\]
The types of $z$ and $\mathsf Tz$ must belong to the domain on which this difference is defined.
\begin{proposition}[Local increment criterion]\label{prop:localincrement}
A decoder $d$ with
\begin{equation}
d(\eta(z))=\delta(z)
\tag{RC5}\label{eq:localincrement}
\end{equation}
exists on the admitted image if and only if $\delta$ is constant on every fiber of $\eta$. Equivalently, $\eta(z)=\eta(z')$ must imply $\delta(z)=\delta(z')$.
\end{proposition}
\begin{proof}
A single decoder value cannot equal two different increments on the same available data. Conversely, constancy on a fiber lets one assign that common increment to its image value; this defines $d$ independently of the chosen preimage.
\end{proof}
For incremental relay, the receiving interface also needs the current scalar or a reference and the prior increments. An increment alone determines only a ratio, not an absolute scalar. Alternatively it may compute $\mathcal C(\mathsf Tz)$ directly from sufficient relational output. The same fiber criterion applies to that target value. Separate scalar transmission is unnecessary only when these available relational data really are sufficient.

A useful failure witness is therefore a pair of admitted inputs with identical available local data but different required scalar updates. This would refute that local relay interface, even if the complete transduction remains injective. The missing information must then be available through an admitted context or receipt interface; it cannot be reconstructed merely by naming it conserved.

\begin{example}[A declared count and a local failure witness]
In Example~7.2, interpret $u,v$ as the presence flags for two distinguished possible relations and choose the illustrative count $\mathcal C(u,v)=u+v$. The invertible map $\mathsf T(u,v)=(u+v\bmod2,v)$ gives
\[
(0,1)\longmapsto(1,1),\quad\Delta\mathcal C=1;
\qquad
(1,0)\longmapsto(1,0),\quad\Delta\mathcal C=0.
\]
Both inputs have complexity one and the same propagated bit $p=1$, yet they require different increments. The propagated bit plus the prior scalar is insufficient; the retained $v$ distinguishes them. The complete map preserves all four input possibilities while this selected structural readout can change. This is a bounded example of the accounting distinction, not a proposed general Thalean complexity measure.
\end{example}


\subsection*{Invariant reconstruction as a concrete special case}
For a supplied represented state $X$ and invariant form $\mathsf G$, an admitted isometry $U$ obeys
\[
I_{\mathsf G}(UX)=(UX)^\dagger\mathsf G(UX)=X^\dagger\mathsf G X.
\]
A receiving interface with sufficient relational state can reconstruct the value without receiving a separate scalar label. This example concerns $I_{\mathsf G}$, not an unproved identification with $\mathcal C$. A receiver holding only one coordinate or one field region generally cannot reconstruct the full invariant; the fiber criterion above is unchanged. The preservation law does not bypass the information needed by the decoder.

\subsection*{Repeated returns and the complexity boundary}
On a domain of positively composed, retained histories, suppose a primitive return $\gamma_x$ and the chosen complexity obey additivity under the specified concatenation. Then induction gives
\begin{equation}
\mathcal C(\gamma_x^{\circ n})=n\mathcal C(\gamma_x),
\qquad n\in\N_0.
\tag{RC6}\label{eq:returncomplexity}
\end{equation}
The base vertex $x$, the repetition count $n$, the history $\gamma_x^{\circ n}$, and any abstract return operator $R_x^n$ remain distinct. A scalar character that is linear in a return exponent does not, by itself, prove that a general relational-complexity functional has that value. Here (RC6) is a conditional specialization, not a new source-native return theorem.

Nor can a nonnegative additive complexity on a full group assign positive cost to an element and remain additive through inverse cancellation: $0=\mathcal C(1)=\mathcal C(g)+\mathcal C(g^{-1})$ would force both terms to vanish. Positive retained-history complexity and signed net return charge therefore require different domains or different laws. Distinct support on a finite carrier can saturate while a retained history grows; neither statistic automatically defines the other.

\subsection*{Complexity normalization is not causal normalization}
A model may select a primitive update with $\Delta\mathcal C=1$, in which case $S_{i+1}=GS_i$. Such a rule needs the chosen functional, its unit, and an admitted update satisfying it. The locked $\cth=1$ concerns field-address reach per event depth and does not imply this complexity increment. Reversible redistribution, unchanged readouts, and admitted uncomputation need not increase complexity. Section~\ref{sec:causalslope} retains its original domain and proof.
\endgroup
\setcounter{section}{11}


\section{Becoming}
Becoming is the ordered composition of admitted realized events. In the new language, it is information-preserving transformation of registered relational organization:
% Original equation number(s) 17; expression unchanged.
\setcounter{equation}{16}
\begin{equation}
\Gamma=e_n\circ\cdots\circ e_2\circ e_1,
\end{equation}

The output context of each event must match the next input. Composition supplies order before any duration in seconds. Its order can change the endpoint and receipt. Algebraic reversal, forward execution under boundary conditions, and thermodynamic irreversibility have separate admission requirements.

\section{Registered History}\label{sec:history}
Registered history retains the distinctions needed to identify a composition relative to its declared projection. One representation is
% Original equation number(s) 18; expression unchanged.
\setcounter{equation}{17}
\begin{equation}
\Gamma=((e_1,r_1^{\rm evt}),\ldots,(e_n,r_n^{\rm evt})),
\end{equation}

where the superscript distinguishes event receipts from the special cycle class $r_1$. A literal list is unnecessary when a compression preserves all distinctions required by the chosen interface. Endpoint permutation, parity, winding, predictive class, and complete ordered word are different such compressions.

Complete contexts may be objects of a category, with admitted histories as arrows. Composition is concatenation modulo explicitly admitted relations. Inverses yield a groupoid only where they are defined and admitted. A finite controller can access an extensible record, but then that record belongs to the complete state. Nonperiodic chronology requires an actual mechanism beyond closed finite deterministic recurrence; a passive log alone supplies no such mechanism.

\begingroup
\renewcommand{\thesection}{13A}
\section{Mode and Phase of Registered Events}\label{sec:mode}
\status{definition on imported native orientation data; finite algebra proved and checked here}
\subsection*{Native input and the common domain}
The native $G_{60}$--$C_5$--$C_3$ orientation mechanism supplies three reflection frames and the inverse order-three orientation pair. It identifies the native designation torsor $\{b,ab\}$ equivariantly with that inverse pair, up to one global orientation reversal \citep[Secs.~8--10]{cave2026orientation}. This is the input used here; it is distinct from the abstract $V_4$ triality of Section~\ref{sec:carrier}.

Use cycle notation on three labels and write
\begin{equation}
\Cevt=(0\ 1\ 2),\qquad \Mevt=(1\ 2),\qquad
\Cevt^3=1,\quad \Mevt^2=1.
\tag{M1}\label{eq:modalgenerators}
\end{equation}
The subscript distinguishes $\Cevt$ from the six-axis conference matrix $C$. Let $\mathcal T$ be the three transpositions in $S_3$. The relative closure-frame set is
\[
\mathscr F=\{(r,\rho)\in\mathcal T^2:r\neq\rho\},\qquad
q(r,\rho)=r^{-1}\rho\in\{\Cevt,\Cevt^{-1}\}.
\]
It has six members. The native correspondence supplies the interpretation of $r$ as closing reflection, $\rho$ as residue, and $q$ as the retained relative orientation. Neither $r$ nor $\rho$ alone determines $q$ over the full relabeling orbit. The action is simultaneous conjugation,
\[
g\cdot(r,\rho)=(grg^{-1},g\rho g^{-1}),
\qquad q(g\cdot f)=gq(f)g^{-1}.
\]

\subsection*{Definition and conjugacy}
\begin{definition}[Mode]\label{def:mode}
Mode is the $S_3$ orientation parity of the registered event cycle. On the registered frame orbit, choose one reference presentation $f_*$ and define
\begin{equation}
\mu(g\cdot f_*)=\sgn(g)\in\{+1,-1\}.
\tag{M2}\label{eq:modedefinition}
\end{equation}
The two modes are the two orientation classes; the reference fixes which class is called positive.
\end{definition}
\begin{proposition}[Registered modal-phase algebra]\label{prop:modealgebra}
The conjugation action on $\mathscr F$ is free and transitive. The native three-cycle and a register-reversing reflection satisfy
\begin{equation}
\Mevt\Cevt\Mevt=\Cevt^{-1},\qquad
\langle\Cevt,\Mevt\rangle\cong S_3.
\tag{M3}\label{eq:modalconjugacy}
\end{equation}
Its six frame presentations split into the two three-element sheets corresponding to
\[
\{1,\Cevt,\Cevt^2\},\qquad
\{\Mevt,\Mevt\Cevt,\Mevt\Cevt^2\}.
\]
\end{proposition}
\begin{proof}
Two distinct transpositions have trivial common centralizer in $S_3$: a permutation commuting with both commutes with their generated group and hence belongs to its trivial center. Thus every frame stabilizer is trivial. Since $|\mathscr F|=|S_3|=6$, each orbit is the full frame set. The expression in \eqref{eq:modedefinition} is therefore well defined.

Conjugating $(0\ 1\ 2)$ by $(1\ 2)$ gives $(0\ 2\ 1)$, its inverse. These permutations generate all six elements of $S_3$. The kernel of the sign character is $A_3=\langle\Cevt\rangle$, so it gives precisely the two sheets. Even relabelings preserve $q$ and odd relabelings exchange its inverse values, in agreement with the imported orientation correspondence.
\end{proof}

\subsection*{Phase, presentation, and event identity}
Each framed presentation has a unique expression $\Cevt^p\Mevt^m\cdot f_*$, with $p\in\Z_3$ and $m\in\Z_2$. Phase is $p$ and Mode is $(-1)^m$. Their action is
\begin{equation}
\Cevt:(p,m)\mapsto(p+1,m),\qquad
\Mevt:(p,m)\mapsto(-p,1-m).
\tag{M4}\label{eq:modalcoordinates}
\end{equation}
These are six coordinates with the semidirect-product law $C_3\rtimes C_2$, not commuting independent actions. The two modes carry the same event cycle with opposite orientation.

The event identity retained by this construction is its declared invariant orbit under the presentation action. This is a statement about registered presentation, not an identification of successive completed occurrences. A general $S_3$ action need not have six-point orbits; the six-point assertion above follows from the explicit free frame action. Absolute plus/minus names depend on the reference, while the two-class structure and the relative sign between presentations do not.

The $b/ab$ pair supplies orientation \emph{states}; $\Mevt$ is an \emph{operation} exchanging them. The relative product $q$ has order three even though its orientation torsor has two members. None of these objects is a binary reduction of the integer time register by definition.

\subsection*{Verification scope}
The accompanying modal check enumerates all six group elements, six frames, and 36 action--frame pairs. It checks conjugacy, freeness, the relative decoder, the parity partition, and the coordinate action. The native provenance is the cited orientation theorem; the present replay verifies its finite frame algebra. A map from this frame action to the $r_1$ time register, the Born coefficient space, or an orientation-sensitive transport must specify its domain and covariance separately.
\endgroup
\setcounter{section}{13}

\begingroup
\renewcommand{\thesection}{13B}
\section{WXYZTI Functionality and the Kernel Crosswalk}\label{sec:wxyzti}
\status{bounded overlay replay; cross-register identification remains a hypothesis}
The recovered WXYZTI data supply four closed six-step register circuits, with role order
\[
IW\to WX\to XY\to YZ\to ZT\to TI\to IW.
\]
At $IW,XY,ZT$, the shared-$B$ move preserves $B$; at $WX,YZ,TI$, the reverse-partner move is $(A,B,C)\mapsto(A,C,B)$. Opposite occurrences belong to the same reciprocal channel:
\begin{equation}
\{IW,YZ\},\qquad \{XY,TI\},\qquad \{ZT,WX\}.
\tag{W1}\label{eq:wxyzti-channels}
\end{equation}
The source notes distinguish this six-step overlay domain from the original ten-step registered word $WXYZTITZYXW$ \citep{wxyztirows,wxyzticontract}. No six-step return result is transferred silently to that ten-step contract.

\begin{proposition}[Bounded reciprocal-role grammar]\label{prop:wxyzti-bounded}
On the supplied four circuits, the $24$ row occurrences split into $12$ shared-$B$ and $12$ reverse-partner rows. Crossing rows of the two types within the same channel gives $48$ candidates. The reciprocal equations
\begin{equation}
S.C_{\rm to}=R.B_{\rm from},\qquad
R.B_{\rm to}=S.C_{\rm from}
\tag{W2}\label{eq:wxyzti-selector}
\end{equation}
select exactly $12$ pairs, four per channel, using each row once. With $j\in\Z_6$ the occurrence position in a fixed circuit, the operations
\[
C_W:j\mapsto j+2,\qquad R_W:j\mapsto j+3
\]
satisfy $C_W^3=R_W^2=1$ and $R_WC_W=C_WR_W$. The latter exchanges reciprocal roles within a channel, not channel orientation.
\end{proposition}
\begin{proof}[Bounded replay]
The script \path{scripts/verify_wxyzti_scope.py} reads the four explicit register rows from the archived source note, reconstructs their consecutive transitions and tests the two move laws. It enumerates all $48$ same-channel candidates and obtains exactly the opposite-edge pairs in the four circuits. The two displayed permutations commute as additions modulo six and are checked on all $24$ occurrences.
\end{proof}
This replays a realized-row universe; it does not derive that universe admission-blind. Some reverse rows were inferred in the early reconstruction and then independently confirmed by the later source packet; the present script is a replay of the archived table, not a claim to reproduce every upstream producer \citep{wxyzticontract}. The reciprocal consistency equations also follow from six-step closure under the stated swap/preservation rules. They cannot be counted twice as independent selection evidence.

\subsection*{Why reciprocal role is not yet Mode}
The modal theorem in Section~\ref{sec:mode} has $\Mevt\Cevt\Mevt=\Cevt^{-1}$. The overlay operation above instead has $R_WC_WR_W=C_W$. This excludes identifying those two \emph{specified} operations under a common action intertwiner. It does not forbid a different native reflection from acting on the WXYZTI domain.

Neither uniform preservation nor uniform flipping of a binary role coordinate alone settles commutation. If $R(p,\mu,r)=(p,\mu,1-r)$ and
\[
M_k(p,\mu,r)=(-p,1-\mu,r+k),\qquad k\in\F_2,
\]
then both $M_0$ and $M_1$ commute with $R$. A noncommuting extension would require an actual different action law, not just the observation that two bits flip together.

\subsection*{The three twelve-element objects}
The twelve reciprocal pairs are indexed by three channels and four realized circuits. Each pair already contains \emph{both} reciprocal roles; the within-pair role bit has been removed by forming the pair. In contrast, $X_{12}$ has three cells and four within-cell points, while the Born object in Section~\ref{sec:born} consists of twelve registered pentagons. Matching cardinalities does not identify any two of these sets.

The proposed link is therefore retained as a source-selection and covariance problem, not a closed kernel theorem:
\begin{hypothesis}[Functional kernel crosswalk]\label{hyp:wxyzti-kernel}
A native registration of WXYZTI occurrences and retained context into points or directed relations of $\mathscr K_{\rm rel}$ preserves the reciprocal coupling and channel continuation. On a separately specified common domain, an admitted reflection realizes the QR modal conjugacy and distinguishes relation direction from binary character.
\end{hypothesis}
A certificate must declare whether its target consists of \emph{points}, \emph{relations}, or a quotient of histories; give a source-derived map and all required additional context; verify every admitted transition and covariance equation; and report collisions or nonuniqueness. Identifying the character $c$ with reciprocal role or the transpose $F_c\leftrightarrow R_c$ with $\Mevt$ is part of this certificate, not an input justified by notation. To connect to the Born operator, an additional map must preserve the signed incidence or Gram structure, rather than merely map twelve labels bijectively.

The existing source reconciliation leaves a native occurrence-dependent endpoint/path assignment open. A fixed assignment of each role letter to one native vertex followed by literal edge reversal cannot implement the required nontrivial return: a path followed by its reverse lifts back to its starting point. The kernel comparison must respect this obstruction rather than relabel it away \citep{wxyzticontract}.
\endgroup
\setcounter{section}{13}

\begingroup
\renewcommand{\thesection}{13C}
\section{Scalar Complexity and Relational Orientation}\label{sec:scalarmode}
\status{descriptive separation and conditional covariance; the native Mode theorem is unchanged}
The intended reading of the scalar is \emph{relational complexity of the thing}. A phase or orientation changes how that structure is presented, not necessarily how much structure it contains. In this sense the two descriptions are orthogonal:
\[
\boxed{\mathcal C(R)}\ \perp_{\rm desc}\
\boxed{\text{phase, Mode, and relational orientation}}.
\]
The subscript specifies descriptive separation, not a supplied inner product, statistical independence, or a proof that all values can be freely combined.

Suppose an admitted presentation action $g$ takes $R$ to an equivalent registered presentation. Definition~4A.2 then gives
\begin{equation}
\mathcal C(g\cdot R)=\mathcal C(R),\qquad S(g\cdot R)=S(R).
\tag{RC7}\label{eq:scalarinvariance}
\end{equation}
When the $S_3$ frame action of Section~\ref{sec:mode} has been attached to such objects on a common domain, a compatible descriptive chart $(c,p,m)$ has
\begin{equation}
\begin{aligned}
\Cevt:(c,p,m)&\longmapsto(c,p+1,m),\\
\Mevt:(c,p,m)&\longmapsto(c,-p,1-m),
\end{aligned}
\qquad p\in\Z_3,\quad m\in\Z_2.
\tag{RC8}\label{eq:complexitymodechart}
\end{equation}
The first coordinate is unchanged; the last two retain exactly the semidirect action (M4). This is a compatible extension under the stated presentation hypothesis, not a newly derived cross-register map.

A reflection can reverse a signed coordinate in a real representation without reversing the nonnegative complexity of its underlying object. In particular, neither $\mathcal C\mapsto-\mathcal C$ nor $S\mapsto-S$ follows from the retained Mode theorem. A signed amplitude, if required, must be separately defined and its covariance proved. Likewise a genuinely structure-changing operation need not preserve $\mathcal C$ merely because it is informally called a reversal.

The motivating sketch placed scalar complexity beside an oriented return about a distinguished null. It helps keep the descriptive roles separate, but no identification of that null with the native phase center is used in a definition or proof here. The null-well, event-frame, and scalar domains require an explicit common-domain construction before the picture can become a native theorem.

\subsection*{When orthogonality becomes geometric}
Descriptive separation does not supply a metric. If an independently declared orientation representation has $\mathbf n\cdot\mathbf n=1$, differentiating gives the actual metric statement $\mathbf n\cdot\dot{\mathbf n}=0$. Section~\ref{sec:gyroscopic} proves the resulting tangent and stabilizer geometry, and states the additional conditions for uniform precession. This metric example neither redefines the capitalized native Mode nor proves that a pure mechanical orientation change preserves an as-yet unconstructed $\mathcal C$.
\endgroup
\setcounter{section}{13}

\begingroup
\renewcommand{\thesection}{13D}
\section{Incidence-Based Orientation Reconciliation}\label{sec:incidencereconciliation}
\status{finite reference and incidence propositions proved here; the native joint history--analyzer domain remains a construction target}
Reference selection and orientation comparison are different operations. A reference names the two sides of a local distinction. Incidence specifies which presentations meet; a comparison law specifies how the orientations at that meeting agree. This section states what can be proved without pretending that all binary labels already belong to one native object.

\begin{definition}[Orientation torsor and reference signing]\label{def:referencetorsor}
A two-member orientation torsor is a set $T$ with a fixed-point-free involution $\iota$. Selecting $t_*\in T$ defines
\begin{equation}
\sigma_{t_*}(t_*)=+1,\qquad
\sigma_{t_*}(\iota t_*)=-1.
\tag{RF4}\label{eq:referencesigning}
\end{equation}
The torsor has an exchange action, but no selected member before the reference is supplied. It is not the acting group $C_2$, whose identity element is already distinguished.
\end{definition}

\begin{proposition}[Change of reference]\label{prop:referencechange}
For $t,t_1,t_2\in T$,
\begin{equation}
\begin{aligned}
\sigma_{\iota t_*}(t)&=-\sigma_{t_*}(t),&
\sigma_{\iota t_*}(\iota t)&=\sigma_{t_*}(t),\\
\sigma_{\iota t_*}(t_1)\sigma_{\iota t_*}(t_2)
&=\sigma_{t_*}(t_1)\sigma_{t_*}(t_2).
\end{aligned}
\tag{RF5}\label{eq:referencechange}
\end{equation}
There are exactly two reference signings. No member of the unmarked torsor is invariant under its exchange automorphism.
\end{proposition}
\begin{proof}
The selected and opposite members exhaust $T$. Reversing the reference interchanges their signs; reversing the represented member at the same time interchanges them again. A product of two coordinates acquires two minus signs. An exchange-fixed selection would require $\iota t=t$, contradicting freeness.
\end{proof}
This proposition does not assert that every richer native presentation has the same automorphism group. An independently supplied marking can reduce that group. It also does not construct a map between two different torsors merely because either can be given $+/-$ coordinates.

\subsection*{Signed incidence as an exact reconciliation mechanism}
Let $E$ be a finite set of edges with selected orientations and $\mathcal H$ a finite family of oriented cycles. In the simple-cycle construction used here, set
\begin{equation}
B_{e\gamma}=\begin{cases}
+1,&\gamma\text{ traverses }e\text{ with its selected orientation},\\
-1,&\gamma\text{ traverses }e\text{ against that orientation},\\
0,&e\text{ is not traversed by }\gamma.
\end{cases}
\qquad M=B^TB.
\tag{IR1}\label{eq:signedincidencereconciliation}
\end{equation}
For more general walks the coefficient is the signed traversal count, not necessarily $0,\pm1$. At a shared simple-cycle incidence, $B_{e\gamma}B_{e\delta}$ records agreement or opposition of two traversals. The sum $M_{\gamma\delta}$ can cancel contributions from different incidences; it is an overlap readout, not the full ordered encounter history.

\begin{proposition}[Incidence reference covariance]\label{prop:incidencecovariance}
Let $D_E,D_H$ be diagonal sign matrices changing the selected edge and history orientations. Then
\begin{equation}
B'=D_EBD_H,\qquad M'=D_HMD_H.
\tag{IR2}\label{eq:incidencegauge}
\end{equation}
If signed permutation actions satisfy
\begin{equation}
U_E(g)B=BU_H(g),
\tag{IR3}\label{eq:incidenceintertwiner}
\end{equation}
then $U_H(g)^TMU_H(g)=M$. Edge-reference changes cancel in $M$; history-reference changes conjugate it and preserve its spectrum.
\end{proposition}
\begin{proof}
Using $D_E^TD_E=I$ gives $B'^TB'=D_HB^TBD_H$. Likewise, transpose and multiply (IR3), using $U_E(g)^TU_E(g)=I$. Diagonal signs and signed permutations are orthogonal, so their conjugations preserve eigenvalues.
\end{proof}
The package-local $G_{15}$ census gives the retained $30\times12$ matrix $B_{\rm hist}$. In this revision, all 120 automorphisms of that local graph are enumerated, their signed edge and cycle actions are constructed, and (IR3) is checked exactly. All 14,400 ordered action products are checked. Omitting the cycle signs fails the intertwining test. This supplies an explicit orientation-reconciliation example on the already admitted incidence construction, not a new identification of these twelve cycles with the Audit014 history alphabet $\{2,3\}$ \citep{cave2026draft6replays}.

\subsection*{The common domain is an incidence relation}
For the open history--analyzer bridge, the proposed architecture is a span
\begin{equation}
\mathcal X_H\ \xleftarrow{\ p_H\ }\ \mathcal I\
\xrightarrow{\ p_A\ }\ \mathcal X_A.
\tag{IR4}\label{eq:incidencespan}
\end{equation}
Here $\mathcal I$ must contain actual admissible joint incidences with the necessary boundary, role, endpoint and phase data. The maps say which history and analyzer presentations participate in each incidence. Neither presentation is required to be the other. Neither map, nor the domain itself, is obtained from equal cardinalities.

Suppose a common group $\Lambda$ acts on this span and lifts to two-member orientation torsors over its endpoint presentations. Pull their signed action coefficients back to $\mathcal I$, writing $c_H(g,\xi)$ and $c_A(g,\xi)$. Both must satisfy the action law
\[
c_j(gh,\xi)=c_j(g,h\xi)c_j(h,\xi),\qquad j\in\{H,A\}.
\]
An incidence comparison is a $C_2$-equivariant bijection $K_\xi$ between the two endpoint torsors. Its coordinate $\kappa(\xi)\in\{\pm1\}$ obeys
\begin{equation}
\sigma_A=\kappa(\xi)\sigma_H,\qquad
\kappa(g\xi)c_H(g,\xi)=c_A(g,\xi)\kappa(\xi).
\tag{IR5}\label{eq:comparisonnaturality}
\end{equation}
The second equation is precisely the commutative-square requirement
$A_gK_\xi=K_{g\xi}H_g$. It is the native comparison law to be derived, not a definition that can be fitted to the desired output after the event.

Under independent changes of local sign convention,
\begin{equation}
\begin{aligned}
\kappa'(\xi)&=\epsilon_A(\xi)\kappa(\xi)\epsilon_H(\xi),\\
c'_j(g,\xi)&=\epsilon_j(g\xi)c_j(g,\xi)\epsilon_j(\xi).
\end{aligned}
\tag{IR6}\label{eq:comparisongauge}
\end{equation}
Substitution preserves (IR5). These are coordinate changes on the pulled-back torsors; changes inherited from endpoint references are the corresponding restricted choices of $\epsilon_j$. Agreement is carried by the incidence and its comparison, not by an obligation to give every local account the same positive name.

\begin{proposition}[Common-domain reconciliation criterion]\label{prop:reconciliationcriterion}
Assume the specified action of $\Lambda$ on a nonempty $\mathcal I$ is transitive and the two coefficients above are honest $C_2$ cocycles. A comparison satisfying (IR5) exists if and only if
\begin{equation}
c_A(h,\xi_0)=c_H(h,\xi_0)
\quad\text{for every }h\in\operatorname{Stab}_{\Lambda}(\xi_0).
\tag{IR7}\label{eq:isotropycriterion}
\end{equation}
When it exists, there are exactly two comparisons, differing by an overall sign. For a finite domain with $N$ action orbits, the criterion applies separately to each orbit; if it holds on all of them, there are $2^N$ comparisons.
\end{proposition}
\begin{proof}
Put $\lambda(g,\xi)=c_A(g,\xi)c_H(g,\xi)$. This is a sign cocycle, and (IR5) becomes $\kappa(g\xi)=\lambda(g,\xi)\kappa(\xi)$. Necessity follows by taking $g=h$ fixing $\xi_0$. For sufficiency, select $\kappa(\xi_0)=\pm1$ and set $\kappa(g\xi_0)=\lambda(g,\xi_0)\kappa(\xi_0)$. If $g\xi_0=g'\xi_0$, then $g=g'h$ for an isotropy element $h$. The cocycle law and (IR7) show that the assigned signs coincide. The same law proves covariance. Transitivity determines every value from the initial one, and the two initial choices give the two solutions. Different orbits impose no equations on one another.
\end{proof}
The theorem removes the demand for an absolute sign without removing the native coupling requirement. A native construction must still supply $\mathcal I,p_H,p_A$, the common action and both lifted actions. It must also establish the relevant orbit structure. ``One remaining global bit'' is justified only after those requirements and the isotropy test are met. On a larger domain there may be several independent signs, or no compatible comparison at all. Selecting a reference cannot repair unequal isotropy characters.

\begin{figure}[htbp]
\centering
\begin{tikzpicture}[>=Latex, every node/.style={font=\small},
box/.style={draw=qrink,rounded corners=0pt,align=center,text width=38mm,minimum height=13mm}]
\node[box] (hist) at (0,0) {History presentation\\and its orientation torsor};
\node[box] (inc) at (5,0) {Admitted incidence $\xi$\\and comparison $K_\xi$};
\node[box] (ana) at (10,0) {Analyzer presentation\\and its orientation torsor};
\draw[->] (inc) -- node[above] {$p_H$} (hist);
\draw[->] (inc) -- node[above] {$p_A$} (ana);
\node[align=center] at (5,-1.25) {Selected references give $\sigma_H,\sigma_A$;\\the incidence law must make $\sigma_A=\kappa(\xi)\sigma_H$ covariant.};
\end{tikzpicture}
\caption{A comparison architecture, not an already derived native joint domain. The bare face-incidence and signed cycle-incidence examples are constructed separately; their combination with the Audit014 history action remains a theorem target.}\label{fig:referenceincidence}
\end{figure}
\endgroup
\setcounter{section}{13}


\partstart{History and Thalean Time}
\section{State Recurrence versus Historical Recurrence}\label{sec:historyconsequences}
For a projection $q$, a returning readout can coexist with a different complete state:
% Original equation number(s) 19; expression unchanged.
\setcounter{equation}{18}
\begin{equation}
q(z_n)=q(z_0),\qquad z_n\neq z_0.
\end{equation}

A two-sheeted cycle with odd holonomy returns to its base after one traversal and to its initial sheet after two. A separately retained integer winding distinguishes every traversal. These are base closure, deck return, and historical winding, respectively. Conversely, an admitted uncomputation can remove a temporary receipt. Record growth is not assumed to be monotone under every reversible operation.

\section{Complete QR State and Predictive History}\label{sec:predictive}
A complete operational state retains the distinctions relevant to future admitted response. For a history $h$, let $\mathcal L(h)$ be its admitted future language and $\mathcal O_h(w)$ its registered response to $w$. Define
% Original equation number(s) 20; expression unchanged.
\setcounter{equation}{19}
\begin{equation}
h_1\sim_{\rm pred}h_2
\quad\Longleftrightarrow\quad
\begin{cases}
\mathcal L(h_1)=\mathcal L(h_2),\\
\mathcal O_{h_1}(w)=\mathcal O_{h_2}(w)\quad\forall w\in\mathcal L(h_1).
\end{cases}
\label{eq:predictive}
\end{equation}

\setcounter{theorem}{0}
\begin{proposition}[Predictive quotient]
If future responses compose consistently under prefixing, predictive equivalence is an equivalence relation and is preserved by appending the same admitted event.
\end{proposition}
\begin{proof}
Reflexivity, symmetry, and transitivity follow from equality of languages and response functions. If $h_1\sim_{\rm pred}h_2$ and $e$ is admitted at both, a future word after $h_i e$ is exactly a word $w$ for which $ew\in\mathcal L(h_i)$. Equality of the original languages and responses gives equality for the residual languages and responses.
\end{proof}

Equality only on common words would make disjoint future languages vacuously indistinguishable. Operational memory therefore requires both admission and response information. An archive can retain more than this predictive quotient; predictive completeness and microscopic reconstructibility are distinct criteria.

\section{Closure}
Closure brings histories into a common comparison context. For histories with common endpoints and admitted reversals, $\gamma_2^{-1}\circ\gamma_1$ is one possible comparison. A remaining deck element, orientation, or history class is its closure residue. Comparison need not reset either history.

Feedback requires an actual dependence on that residue:
% Original equation number(s) 21; expression unchanged.
\setcounter{equation}{20}
\begin{equation}
\Adm(z,r)\neq\Adm(z,r')
\quad\text{or}\quad
\mathcal O_{z,r}(w)\neq\mathcal O_{z,r'}(w).
\end{equation}

Distinct residue labels alone do not establish either inequality. A source or update law must provide the coupling; the comparison operation itself asserts no optimization principle.

\begingroup
\renewcommand{\thesection}{16A}
\section{Literal History Phase Transport and Its Scope}\label{sec:historyphasegroupoid}
\status{exact replay from supplied native permutations and explicitly transcribed source indices; no decorated history-state identification inferred}
The registered-history phase object provides a concrete test of what reference selection can and cannot settle. To avoid collision with the correlation $\Gamma(A,B)$ or a history denoted $\Gamma$, write its group as $\Gamma_{\rm hist}$. The author-supplied Project~41 excerpts specify a sixteen-element subgroup of the cached $\Aut(G_{60})$ action with
\begin{equation}
\begin{gathered}
\Gamma_{\rm hist}=\langle r,s,z\rangle\cong D_8\times C_2,\\
r^4=s^2=z^2=e,\qquad srs=r^{-1},\qquad [z,r]=[z,s]=e,\\
r=\Aut(G_{60})[3],\quad s=\Aut(G_{60})[325],\quad z=\Aut(G_{60})[324].
\end{gathered}
\tag{HG1}\label{eq:historygrouppresentation}
\end{equation}
The literal cache, rather than a group fingerprint, is used in the new replay. It verifies the presentation, the internal direct product, and generation by $r$ together with the source stabilizer. Completeness of the supplied 480-element automorphism census remains imported; permutation composition and graph preservation are checked directly \citep{cave2026nativepacket,cave2026p4receipts,cave2026draft6replays}.

Set $a=r^2z$, $H_0=\langle a,sz\rangle$, and $K=\langle a\rangle$. The phase set is the left coset space $X=\Gamma_{\rm hist}/H_0$. Its four members and isotropy are
\begin{equation}
S_k=r^kH_0,\quad H_k=r^kH_0r^{-k},\quad
w_k=r^{2(k\bmod2)}sz,\quad H_k=\langle a,w_k\rangle.
\tag{HG2}\label{eq:historyphasecosets}
\end{equation}
The computed $r$ powers have native indices $[0,3,2,5]$, $r^{-1}$ has index $5$, and $K=\{0,326\}$. The native indices are coordinates in the supplied ordered cache, not universal names.

\begin{center}
\small
\begin{tabular}{@{}ccll@{}}
\toprule
$k$ & $r^k$ index & $H_k$ indices & $w_k$ index \\
\midrule
0 & 0 & $\{0,1,323,326\}$ & 1 \\
1 & 3 & $\{0,7,325,326\}$ & 7 \\
2 & 2 & $\{0,1,323,326\}$ & 1 \\
3 & 5 & $\{0,7,325,326\}$ & 7 \\
\bottomrule
\end{tabular}
\end{center}
The phase has period four, while this isotropy and word-generator family has period two. Every $H_k/K$ is a group of order two. Its nonidentity element acts as a history-word exchange in the source interpretation; the quotient group itself is not a torsor with an unnamed identity.

\begin{proposition}[Phase-transport action-groupoid automorphism]\label{prop:phasegroupoid}
For the declared action groupoid $\Gamma_{\rm hist}\ltimes X$, define
\begin{equation}
\alpha(g)=rgr^{-1},\qquad \tau(x)=rx.
\qquad \tau(gx)=\alpha(g)\tau(x).
\tag{HG3}\label{eq:phasegroupoidautomorphism}
\end{equation}
Then $\tau$ on objects and $g\mapsto\alpha(g)$ on arrow labels define an automorphism. It transports $H_k$ to $H_{k+1}$, $S_k$ to $S_{k+1}$ and $w_k$ to $w_{k+1}$, and fixes $K$ setwise.
\end{proposition}
\begin{proof}
Conjugation is a group automorphism, and $r(gx)=(rgr^{-1})(rx)$. Thus an arrow $g:x\to gx$ is sent to $\alpha(g):\tau(x)\to\tau(gx)$. Identities, inverse arrows and composition are preserved; replacing $r$ by $r^{-1}$ gives the inverse functor. Formula (HG2) gives the isotropy and coset identities. The supplied presentation gives $r w_k r^{-1}=w_{k+1}$ and centrality of $a$. All identities are also checked on the literal indexed permutations and all four objects.
\end{proof}
This is the exact finite content of the phase-transport part of Audit018K, sealed in the author's receipt at commit \sourcefile{42ea48f}. Selecting $S_0$ as part of the reference makes $r$ a map to another selected reference, not an automorphism fixing that marked object. Reference covariance is not the assertion that a marking remains pointwise fixed.

\subsection*{Actor transport is not yet a history-state crosswalk}
The reduced continuation interface has two labels $h\in\{2,3\}$, with preserve-roles action $h\mapsto5-h$ and swap-roles action $h\mapsto h$. Their state exchange commutes with both reduced actions. Audit018I records exactly two abstract equivariant bijections to a sign torsor, differing by an overall minus sign. These finite facts are replayed separately.

The stronger inference that (HG3) induces the literal $2\leftrightarrow3$ action requires an additional construction. One needs the appropriate decorated history carrier $\mathcal Y$, its projection $\pi_{\rm hist}:\mathcal Y\to\{2,3\}$ and a lifted automorphism satisfying
\begin{equation}
\pi_{\rm hist}\circ\widetilde\tau=\iota\circ\pi_{\rm hist},
\qquad \iota(2)=3,\quad\iota(3)=2.
\tag{HG4}\label{eq:decoratedhistorytarget}
\end{equation}
The shown Audit018K producer combines the phase theorem with the separately imported reduced-swap result; it does not construct this projection and check (HG4). Conjugating one word-generator index to another transports an actor between fibers. It does not, solely by that computation, determine the permutation of the states in those fibers.

Accordingly this manuscript retains the sealed permutation result, qualifies its stronger state-exchange interpretation, and claims no exhaustive no-go for all richer native history presentations. It does not select a native history-to-central-sign crosswalk or a $\Gamma(A,B)$-sign actor. Program~3 VN03/VN04 remain open at the common-domain action interface. A selected reference can remove a naming ambiguity, but only a construction such as (IR4)--(IR7), or an equivalent native comparison, can establish the coupling.
\endgroup
\setcounter{section}{16}


\section{Thalean Time}\label{sec:time}
\status{locked definition on an imported certified registered-history cycle domain}
The Project~41 synthesis reports
% Original equation number(s) 22; expression unchanged.
\setcounter{equation}{21}
\begin{equation}
\rank H_1(\Hreg;\Z)=3,
\qquad \{r_0,r_1,r_2\}
\end{equation}

as its primitive residual register, with exactly $r_1$ visible under the certified $G_{30}$ registration. Its admitted order-four action and extractor satisfy
% Original equation number(s) 23; expression unchanged.
\setcounter{equation}{22}
\begin{equation}
g^2\neq1,\qquad g^4=1,\qquad D_a(g^2)=r_1.
\end{equation}

The native register and selector are imported from Audits 201AE, 201BY, 201AJ, 201CK, and definition record 201CL \citep[Sec.~17]{cave2026becoming}.

\setcounter{theorem}{0}
\begin{definition}[Thalean time]
On the admitted closed registered-history domain, Thalean time is the primitive additive oriented winding of the distinguished $G_{30}$-visible registered-history class $r_1$:
% Original equation number(s) 24; expression unchanged.
\setcounter{equation}{23}
\begin{equation}
\tth(\gamma)=\wind_{r_1}(\gamma)\in\Z.
\label{eq:thaleantime}
\end{equation}

\end{definition}

The declared register supplies the dual functional
% Original equation number(s) 25; expression unchanged.
\setcounter{equation}{24}
\begin{equation}
\varphi_1(r_0)=0,\qquad\varphi_1(r_1)=1,\qquad\varphi_1(r_2)=0.
\end{equation}

so $\tth(\gamma)=\varphi_1([\gamma])$. A primitive class in an arbitrary rank-three group would not by itself determine this dual functional; the visibility normalization and complementary register are part of the input.

\setcounter{theorem}{1}
\begin{proposition}[Temporal composition]
For composable admitted closed histories,
% Original equation number(s) 26; expression unchanged.
\setcounter{equation}{25}
\begin{equation}
\tth(\gamma_2\circ\gamma_1)=\tth(\gamma_2)+\tth(\gamma_1),
\qquad
\tth(\gamma^{-1})=-\tth(\gamma).
\end{equation}

\end{proposition}
\begin{proof}
Homology sends concatenation to addition and reversal to negation. Applying the homomorphism $\varphi_1$ gives both identities.
\end{proof}

An open-history extension needs reference paths, a covering lift, or an extending edge cocycle. Such choices can add endpoint terms while preserving closed-cycle winding. The transport domain used in $\cth$ is therefore not silently identified with the certified closed-history domain.

Mode is a binary orientation class in the registered $S_3$ action. Thalean time is a $\Z$-valued functional on the specified cycle register. The conjugacy of $\Cevt$ does not by itself identify its orientation sign with $\wind_{r_1}$; a common-domain intertwiner would supply that additional relation. The definitions remain separate.

\section{Event Step and Thalean Tick}\label{sec:tick}
An event step is one irreducible operation of the admitted propagation grammar. Grouping several handoffs into a macro-operation does not reduce its cost. A Thalean tick is one positively oriented primitive traversal of $r_1$:
% Original equation number(s) 27; expression unchanged.
\setcounter{equation}{26}
\begin{equation}
\Delta\tth=+1.
\end{equation}

If that traversal requires $d_1$ events, its rate is $1/d_1$ tick per event step. A forward traversal followed by its inverse has zero net winding but positive execution depth. A clock calibration is an additional correspondence, not a replacement of either native quantity.

\section{Event Depth, Recurrence, Frequency, and Temperature}
Recurrence is return under a specified update or readout. A discrete unitary eigenvector $Uv=e^{-i\theta}v$ has eigenphase $\theta$ per update; only after assigning $\Delta t$ does it yield $\omega=(\theta+2\pi k)/\Delta t$. The original four-way register remains
% Original equation number(s) 28; expression unchanged.
\setcounter{equation}{27}
\begin{equation}
\begin{array}{rcl}
\cth&:&\text{causal reach per irreducible event depth},\\
\tth&:&\text{oriented registered-history winding},\\
\omega&:&\text{recurrence or eigenphase rate in a declared parameter},\\
\Theta&:&\text{equilibrium occupation parameter, when justified}.
\end{array}
\end{equation}

Mode adds a separate binary character, while finite event phase is a $C_3$ position. Collective \emph{eigenmodes} and their occupations belong to a chosen dynamics. An arbitrary occupation distribution need not define a temperature.
